\section{Experimental protocols}
\label{app:experiments}

All paths below are relative to the repository root, and all numerical kernels
use float64 arrays. The objective reported throughout is the Gaussian
variational objective, up to an additive target-dependent constant,
\begin{equation*}
  \calE(m,C) = -\frac12\log\det C + \E_{\Theta\sim\N(m,C)}V(\Theta),
\end{equation*}
and $\DeltaE(a) = \calE(a) - \calE(a_\star)$ is the energy gap against the
optimizer of the same experiment. Where deterministic quadrature is used,
NumPy's $K$-point Gauss--Hermite nodes $(z_j,w_j)$ are rescaled by
$\E f(Z) = \pi^{-1/2}\sum_{j\leq K}w_jf(\sqrt2z_j)$ for $Z\sim\N(0,1)$. Writing
$G_n = \E_{\rho_{a_n}}\nabla V$ and $A_n = \E_{\rho_{a_n}}\nabla^2V$, the two
Fisher--Rao steps are \eqref{upd:Riem} and \eqref{upd:KL} in the equivalent
form $m_{n+1} = m_n - \Delta t\,C_nG_n$ with
$C_{n+1}^{\Riem} = C_n^{1/2}\exp[\Delta t\{\Id-C_n^{1/2}A_nC_n^{1/2}\}]C_n^{1/2}$
and $C_{n+1}^{\KL} = (1+\Delta t)(C_n^{-1}+\Delta tA_n)^{-1}$.

Unless stated otherwise a run is evolved for its fixed iteration budget, and a
time to tolerance is a diagnostic computed from the saved trajectory rather
than an early-stopping rule. The PDF files displayed in \Cref{sec:numerics} are
byte-for-byte copies of the corresponding files in
\path{reports/assets/figs}; the script \path{reports/make_report_assets.py}
performs only post-processing and never reruns a dynamics experiment.

Each experiment below belongs to a named group whose configurations live under
\path{configs/<group>}, whose runners live under \path{scripts/<group>}, and
whose kernels and plot selectors live under \path{src/<group>}. The lineage
paragraphs name the group, its configuration and runner files, and, for each
selected panel, the CSV it reads, the builder function that consumes it, and
the resulting file in \path{reports/assets/figs}.

\subsection{Covariance bootstrap}

Let $H = \diag(\nu_1,\ldots,\nu_{N_\theta})$ with
$(\nu_i)_{i=1}^{N_\theta}$ log-spaced in $[1,\kappa]$. The
Gaussian target is $V(\theta) = \frac12\theta^\top H\theta$ with
$(m_\star,C_\star) = (0,H^{-1})$, and the non-Gaussian target used only in the
transport experiment is
$V(\theta) = \frac12\theta^\top H\theta
+ \gamma\sum_{i=1}^{N_\theta}\log\cosh(\theta_i)$ with $\gamma = 1$. For the
latter $m_\star = 0$ and each diagonal entry of $C_\star$ solves
\begin{equation}\label{eq:app-diag-fixed-point}
  c_i^{-1} = \nu_i + \gamma\,\E_{Z\sim\N(0,c_i)}\operatorname{sech}^2Z,
\end{equation}
computed by fixed-point iteration with tolerance $10^{-13}$, at most $500$
iterations, and $80$-point Gauss--Hermite quadrature. Thus $\alpha = 1$, while
$\beta = \kappa$ for the Gaussian target and $\beta = \kappa+1$ for the
$\log\cosh$ target. Every run starts from
$m_0 = m_\star + 3\sqrt{\diag(C_\star)}$ and $C_0 = \lambda_{0,\min}\Id$, and
the Gaussian scaling grid is
\begin{equation*}
  N_\theta\in\{10,20\},\quad
  \kappa\in\{10,100,1000\},\quad
  \lambda_{0,\min}\in\{10^{-10},10^{-8},10^{-6},10^{-4},10^{-2}\},
\end{equation*}
for both covariance maps, with $\Delta t = 0.5$ and $200$ steps. The burn-in
time is the first $n$ such that $\lambda_{\min}(C_n)\geq(2\beta)^{-1}$. The left
panel of \Cref{fig:num-three-stage} selects KL, $N_\theta = 20$,
$\kappa = 1000$, and $\lambda_{0,\min} = 10^{-10}$; its three ordinates are
$\lambda_{\min}(C_n)$, the exact Gaussian energy gap, and
$\norm{C_n^{1/2}G_n}$, and its vertical markers are the covariance threshold
and the first saved iterate with energy gap at most $10^{-3}$.

For the transport grid, the Gaussian and $\log\cosh$ targets both have
$N_\theta = 20$ and $\kappa = 50$, and the grid is
\begin{equation*}
  \lambda_{0,\min}\in\{10^{-8},10^{-6},10^{-4}\},\qquad
  c\in\{0.25,0.5,1\},\qquad \Delta t = 0.5,
\end{equation*}
with both Fisher--Rao tails and a $6000$-step budget. When
$\lambda_{0,\min} < c/\beta$, one Bures step of size $\eta = c/\beta$ is applied
first,
\begin{align*}
  m^+ &= m - \eta G, \qquad \widetilde C = (\Id-\eta A)C(\Id-\eta A),\\
  C^+ &= \frac12\left\{\widetilde C + 2\eta\Id
    + [\widetilde C(\widetilde C+4\eta\Id)]^{1/2}\right\}.
\end{align*}
The right panel of \Cref{fig:num-three-stage} selects $c = 0.5$ and plots, for
every target, tail, and $\lambda_{0,\min}$, the first iteration whose gap is at
most $10^{-1}$, together with $\lambda_{\min}(C)$ after the first step.

\paragraph{Data and figure lineage.}
\begingroup\small\raggedright
The group is \path{natural_gradient_covariance_bootstrap}, with configuration
\path{covariance_bootstrap.yaml} and runner \path{run_experiments.py}. The left
panel reads \path{outputs/natural_gradient_covariance_bootstrap/results_long.csv}
through \path{plotting.py::fig_three_stage} into
\path{fig_covboot_three_stage.pdf}; the right panel reads
\path{outputs/natural_gradient_covariance_bootstrap/wasserstein_bootstrap_summary.csv}
through \path{plotting.py::fig_wasserstein_bootstrap} into
\path{fig_covboot_wasserstein_bootstrap.pdf}.
\par\endgroup

\subsection{Deterministic discretization grid}

All three targets have $N_\theta = 2$, use $\lambda\in\{0.01,0.1,1\}$, and are
\begin{align*}
  V_{\rm G}(\theta) &= \frac12\theta^\top\diag(1,\lambda)\theta,\\
  V_{\rm q}(\theta) &= \frac{(\sqrt\lambda\,\theta_1-\theta_2)^2}{20}
    + \frac{\theta_2^4}{20},\\
  V_{\rm s}(\theta) &= \frac{(\sqrt\lambda\,\theta_1-\theta_2)^2}{20}
    + \frac{\delta}{2}\norm{\theta}^2 + \gamma\log\cosh\theta_2,
  \qquad (\delta,\gamma) = (0.05,1).
\end{align*}
The Gaussian initialization is $m_0 = (10,10)$, $C_0 = \diag(0.5,2)$; the other
two use $m_0 = (10,10)$, $C_0 = 4\Id$. Gaussian and quartic expectations are
analytic through order four, and the smooth target uses $80$-point
one-dimensional Gauss--Hermite quadrature for the $\theta_2$-marginal. The
Gaussian optimum is analytic; the other two optima are computed in Cholesky
coordinates by four deterministic L-BFGS-B starts with a $5000$-iteration
budget, gradient tolerance $10^{-10}$, and function tolerance $10^{-16}$.

The common grid is
$\Delta t\in\{0.001,0.002,0.005,0.01,0.02,0.05,0.1,0.2,0.5,1,2,5,10\}$ with
$T = 15$ and $N = \lceil T/\Delta t\rceil$, and both schemes use exact
population $G_n$ and $A_n$. A Radau solution of \eqref{ODEs:NG}, with relative
tolerance $10^{-10}$, absolute tolerance $10^{-12}$, and $400$ output times, is
used only for the terminal discretization-accuracy diagnostic. Simulation stops
early only on a nonfinite update or $\lambda_{\min}(C_n)\leq10^{-12}$. A
completed trajectory is called stable when it is SPD-feasible, ends below its
initial energy, and never exceeds $10^3$ times its initial energy gap;
monotonicity allows an absolute per-step tolerance of $10^{-10}$; and accuracy
means
$\norm{m_N-m(T)} +
\norm{C_N-C(T)}_{\Frob}/\norm{C(T)}_{\Frob}\leq10^{-2}$.

The left panel of \Cref{fig:num-discretization} uses only $V_{\rm s}$, all three
values of $\lambda$, and the matched stable steps
$\Delta t\in\{0.05,0.1,0.5,1\}$, plotting the energy gap against $n\Delta t$.
The right panel retains every stable run that reaches gap $10^{-6}$, and for
each target, method, and $\Delta t$ averages the hitting iteration and the
hitting time $n\Delta t$ over the eligible values of $\lambda$; its top row is
iteration count and its bottom row is physical flow time, not measured wall
time.

\paragraph{Data and figure lineage.}
\begingroup\small\raggedright
The group is \path{natural_gradient_discretization_stepsize}, with configuration
\path{stepsize_grid.yaml} and runner \path{run_stepsize_grid.py}; the target,
method, optimizer, ODE, and classification code sits in its \path{src} tree.
Both selected panels read
\path{outputs/natural_gradient_discretization_stepsize/results_long.csv} and/or
\path{outputs/natural_gradient_discretization_stepsize/summary.csv}, through the
builders \path{scripts/natural_gradient_discretization_stepsize/plot_results.py::fig_convergence_speed_grid}
and \path{plot_results.py::fig_iterations_vs_stepsize}, into
\path{fig_convergence_speed_smooth_logconcave.pdf} and
\path{fig_iterations_vs_stepsize.pdf}.
\par\endgroup

\subsection{Sticking-the-landing variance and constant-step floors}

Here $A = \diag(a_1,\ldots,a_{N_\theta})$ with $(a_i)$ log-spaced in
$[1,\kappa]$, and the targets are
$V_{\rm G}(\theta) = \frac12\theta^\top A\theta$ and
$V_{\rm lc}(\theta) = \frac12\theta^\top A\theta
+ \tau\sum_{i=1}^{N_\theta}\log\cosh(\theta_i)$. The grid is
$N_\theta\in\{2,8,32\}$ and $\kappa\in\{10,100,10^4\}$ for $V_{\rm G}$, and
$N_\theta\in\{2,8,32\}$, $\kappa\in\{10,100\}$, $\tau\in\{0.1,1\}$ for
$V_{\rm lc}$. For $V_{\rm lc}$ one has $m_\star = 0$ and $C_\star$ diagonal with
$(C_\star)_{ii}$ solving \eqref{eq:app-diag-fixed-point} with $\nu_i$ replaced
by $a_i$ and $\gamma$ by $\tau$, under the same quadrature, tolerance, and
iteration cap.

The estimator-level experiment compares the baseline score
$b_{\rm base}(\theta) = \nabla\log\rho_\post(\theta)$ with the
sticking-the-landing estimator
$b_{\STL}(\theta) = b_{\rm base}(\theta) + C^{-1}(\theta-m)$.
At each target configuration the four states plotted on the left of
\Cref{fig:num-stl} have $C = C_\star$ and
$m = m_\star + \rho\sqrt{\diag(C_\star)}$ with $\rho\in\{0,0.25,1,4\}$; the full
estimator grid additionally includes $(m_\star,0.25C_\star)$ and
$(m_\star,4C_\star)$. For every state and each seed $s\in\{0,\ldots,19\}$,
$M = 100000$ Gaussian samples are drawn in one chunk, the configured chunk cap
being $200000$. The reported mean-block Fisher--Rao variance is
$\Var_{\FR}(b) = \E[(b-\E b)^\top C(b-\E b)]$. The left panel plots, against
$\rho$, the median over the $20$ seeds of $\Var_{\FR}(b_{\STL})$ divided by
$\Var_{\FR}(b_{\rm base})$, separately for
all $12$ non-Gaussian configurations; for display on a logarithmic axis the
plotting code adds $10^{-3}$ to $\rho$ and clips ratios below $10^{-16}$.

The algorithm-level experiment compares
$\{\Riem,\Riem\text{-}\STL,\KL,\KL\text{-}\STL\}$. At every step it draws $B$
samples from the current Gaussian, averages the corresponding mean estimator,
and averages the diagonal pointwise log-target Hessian for the covariance step,
the two mean estimators within a covariance scheme using the same
standard-normal stream. The settings are $\Delta t = 0.1$, $B\in\{1,4,16\}$,
$m_0 = m_\star + 3\sqrt{\diag(C_\star)}$, and $C_0 = 0.5C_\star$,
with $2000$ steps for $V_{\rm G}$ and $3000$ for $V_{\rm lc}$. Gaussian energy
gaps are analytic, and non-Gaussian gaps use $40$-point Gauss--Hermite
quadrature and are therefore independent of the mini-batch used for the update.
An emergency eigenvalue guard clips values below $10^{-12}$; the runner counts
such clips in a cell diagnostic printed during execution, but the surviving
summary CSV does not preserve that count.

The replicate labels are $0,\ldots,19$. For the algorithm grid, one batched
random stream advances these $20$ leading-array entries, and its cell seed is
the first $15$ hexadecimal digits of
$\operatorname{SHA256}(\texttt{20260623|kind|d|kappa|tau|B})$, interpreted as an
integer; the method name is deliberately absent, so all four paired methods
receive the same standard-normal sequence. For the estimator grid, the generator
seed is the replicate label itself.

For each run the empirical floor is the median gap over the final $20$ percent
of the full, undecimated trajectory. The right panel of \Cref{fig:num-stl} takes
the median of these per-run medians over every target configuration and seed,
separately by target kind, method, and $B$, and clips nonpositive display values
to $10^{-16}$.

\paragraph{Data and figure lineage.}
\begingroup\small\raggedright
The group is \path{natural_gradient_stl_variance}, with configuration
\path{stl_variance.yaml} and the two runners
\path{run_estimator_variance.py} and \path{run_algorithm_grid.py}. The left
panel reads
\path{outputs/natural_gradient_stl_variance/estimator_variance.csv} through
\path{plotting.py::fig_variance_by_distance} with \texttt{kind=log\_cosh} into
\path{fig_stl_variance_by_distance_to_optimum.pdf}; the right panel reads
\path{outputs/natural_gradient_stl_variance/tail_noise_floor.csv} through
\path{plotting.py::fig_batchsize_effect} into
\path{fig_stl_batchsize_effect.pdf}.
\par\endgroup

\subsection{Trace and traceless affine modes}

At a state $(m,C)$, set $A_C = \E_{\rho_a}\nabla^2V$ and $B = C^{1/2}A_CC^{1/2}$.
Both experiments run
\begin{equation}\label{eq:app-affine-step}
  m^+ = m - \Delta t\,C\,\E_{\rho_a}\nabla V,
  \qquad
  C^+ = C^{1/2}\exp\!\left[\frac{\Delta t}{2\omega}
    \left\{-B + \frac{\omega+\tau\Tr B}{\omega+N_\theta\tau}\Id\right\}\right]
    C^{1/2}.
\end{equation}
For the Gaussian calibration target $\rho_\post = \N(0,\Id_{N_\theta})$
expectations are exact, so $B = C$, the mean step is $m^+ = m-\Delta t\,Cm$, and
the covariance step diagonalizes as
$\lambda_i^+ = \lambda_i\exp[\frac{\Delta t}{2\omega}\{-\lambda_i+q_\tau(C)\}]$
with $q_\tau(C) = (\omega+\tau\Tr C)/(\omega+N_\theta\tau)$. Its grid is
$N_\theta\in\{2,5,10\}$, $\omega\in\{0.125,0.25,0.5,1,2\}$, and
$\tau\in\{-\omega/(2N_\theta),0,\omega/(2N_\theta)\}$, with $\Delta t = 0.02$,
$T = 20$, and metrics saved every five steps.

The non-Gaussian target is
\begin{equation*}
  V_\rho(\theta) = \frac12\norm{\theta}^2
  + \frac{\rho}{M}\sum_{\ell=1}^{M}\log\cosh(a_\ell^\top\theta),
  \qquad N_\theta = 5,\quad \rho = 5,\quad M = 4N_\theta = 20,
\end{equation*}
where the unit vectors $a_\ell$ are normalized Gaussian draws from NumPy seed
$123$. The expectations in \eqref{eq:app-affine-step} use one fixed scrambled
Sobol design, transformed coordinatewise by the Gaussian inverse CDF, with
$K = 4096$ and seed $2026$, and the same points are reused at every state and
for every $(\omega,\tau)$ and initialization. The reference optimum uses an
independent scrambled Sobol design with $K_{\rm ref} = 8192$ and seed $2027$,
followed by Cholesky-coordinate L-BFGS-B with a $2000$-iteration budget,
gradient tolerance $10^{-6}$, and function tolerance $10^{-15}$. The dynamics
grid is $\omega\in\{0.25,0.5,1\}$ and
$\tau\in\{-\omega/(2N_\theta),0,\omega/(2N_\theta)\}$, with $\Delta t = 0.005$,
$T = 40$, and metrics saved every ten steps. There are no random replicates:
once the seeds are fixed, the quasi-Monte Carlo design is a single common
deterministic one.

The five Gaussian initializations are
\begin{align*}
  \text{mean:}&\quad m_0 = 3\one/\sqrt{N_\theta},\quad C_0 = \Id,\\
  \text{volume-high/low:}&\quad m_0 = 0,\quad
    C_0 = 4\Id\ \text{or}\ \Id/4,\\
  \text{shape:}&\quad m_0 = 0,\quad
    C_0 = \diag(e^2,e^{-2},1,\ldots,1),\\
  \text{mixed:}&\quad m_0 = 2\one/\sqrt{N_\theta},\quad
    C_0 = 2\diag(e^{1.5},e^{-1.5},1,\ldots,1),
\end{align*}
and for the non-Gaussian target the same formulas are applied in the
$(m_\star,C_\star)$-whitened coordinates, so that for instance the shape
covariance is $C_\star^{1/2}\diag(e^2,e^{-2},1,\ldots,1)C_\star^{1/2}$.

Both panels of \Cref{fig:num-affine} show $T(\tau)/T(0)$ for every $\omega$ and
initialization, which is the ratio predicted mode by mode in
\Cref{thm:aff-modes}. The Gaussian panel selects $N_\theta = 5$ and the first
saved time at which the exact Gaussian KL, normalized by its initial value, is
at most $10^{-4}$; the non-Gaussian panel uses the fixed quasi-Monte Carlo
objective gap, normalized by its initial gap, at tolerance $10^{-2}$. A cell is
left blank if either hitting time is not finite within the horizon.

\paragraph{Data and figure lineage.}
\begingroup\small\raggedright
The group is \path{omega_tau_modes}, with configurations
\path{gaussian_target.yaml} and \path{logconcave_target.yaml}. The report
builder reads the legacy committed locations
\path{outputs/gaussian_grid/summary.csv} and
\path{outputs/logconcave_grid/summary.csv} through
\path{reports/make_report_assets.py::_fig_tau_speedup}, producing
\path{fig_ot_gaussian_tau_speedup.pdf} and
\path{fig_ot_logconcave_tau_speedup.pdf}. To reproduce these exact locations
with the current runners, their \texttt{--outdir} arguments must override the
nested \texttt{output\_dir} values in the YAML files.
\par\endgroup

\subsection{Non-log-concave boundary}

All selected runs are centered, scalar, and deterministic, and use $160$-point
Gauss--Hermite quadrature for the double well \eqref{eq:double-well}, whose
Hessian $V_R''(\theta) = -1+2\tanh^2(\theta/R)$ lies in $[-1,1]$. For
$\Theta\sim\N(0,c)$ we write $A_R(c) = \E V_R''(\Theta)$, and the plotted
squared Fisher--Rao gradient is $(1-cA_R(c))^2$.

The Riemannian cascade runs the map \eqref{eq:cascade-map} of
\Cref{prop:cascade-Riem} with $R\in\{20,50,100,300,1000\}$, $c_0 = 1$,
$\Delta t = 1$, and a $20$-step cap, stopping before the next update if
$c_n>R^2/4$, if $\log c$ would exceed $700$, or on nonfinite arithmetic. The
first panel of \Cref{fig:num-nonconvex} plots every recorded iterate for all
five values of $R$.

The pole experiment runs the unprojected resolvent \eqref{eq:KL-cov-map} of
\Cref{prop:pole-KL} with $R = 1000$, $\Delta t = 1$, $c_0 = 1-\varepsilon$,
$\varepsilon\in\{10^{-1},10^{-2},10^{-3},10^{-4},10^{-5},10^{-6}\}$, and a
two-step cap. The middle panel plots the squared gradient after the first step
against $\varepsilon$, together with the reference $4\varepsilon^{-2}$. At
$R=1000$ the sufficient condition $R_\varepsilon^2\geq2\varepsilon^{-2}$ used in
the proof of \Cref{prop:pole-KL} holds only for
$\varepsilon\in\{10^{-1},10^{-2}\}$; the measured-to-reference ratios are
$0.90$ and $0.99$ there, remain $0.99$ and $0.89$ at the uncertified points
$10^{-3}$ and $10^{-4}$, and collapse to $0.39$ and $0.014$ at $10^{-5}$ and
$10^{-6}$, where the fixed truncation radius saturates the pole. The
small-$\varepsilon$ bend of the middle panel is this fixed-$R$ saturation and
not a failure of the $\varepsilon^{-2}$ law along the proposition's admissible
sequences $R_\varepsilon\to\infty$.

The clipped experiment runs the scalar form of \eqref{upd:clip-KL},
\begin{equation*}
  c_{n+1} = \operatorname{clip}_{[0.5,2]}
  \frac{(1+\Delta t)c_n}{1+\Delta t\,c_nA_R(c_n)},
  \qquad c_0 = 1,\qquad \beta = 1.
\end{equation*}
The report panel selects the primary safety fraction $0.5$, hence
$L_{\rm clip} = 2\beta\lambda_+ = 4$ and
$\Delta t = 0.5/L_{\rm clip} = 0.125$, and runs $300$ steps for each of the five
values of $R$. The full output also contains safety fractions $0.25$ and $0.9$,
together with the deliberately outside-theorem step
$\Delta t = 1/(\beta\lambda_+) = 0.5$; the selected panel excludes those rows.
At step $n$ the code records the centered covariance term of the split
stationarity measure \eqref{eq:split-stationarity},
\begin{equation*}
  D_n = \KL(\N(0,c_n)\|\N(0,c_{n+1})),
  \qquad
  D_{\min}(N) = \min_{0\leq n<N}D_n,
  \qquad
  B_N = \frac{\Delta t}{N}\{\calE(0,c_0)-\calE(0,c_N)\},
\end{equation*}
and the right panel plots $D_{\min}(N)$ against the envelope $B_N$ of
\Cref{thm:clip-KL}. In the recorded primary runs the upper clip first activates
at step $2$ for every $R$; the step-$2$ displacement is nonzero, while the next
and subsequent pinned steps have $D_n = 0$.

\paragraph{Data and figure lineage.}
\begingroup\small\raggedright
The group is \path{natural_gradient_nonconvex_instability}, with configuration
\path{nonconvex_instability.yaml}. Writing \path{outputs} for
\path{outputs/natural_gradient_nonconvex_instability}, the cascade panel reads
\path{outputs/results_long.csv} through
\path{plotting.py::fig_riemannian_cascade}, the pole panel reads
\path{outputs/kl_pole_summary.csv} through \path{plotting.py::fig_kl_pole}, and
the clipped panel reads \path{outputs/clipped_kl_stationarity.csv} through
\path{plotting.py::fig_clipped_kl_stationarity}, restricted to
\texttt{dt\_rule=theorem\_safe}. The assets are, respectively,
\path{fig_nonconvex_riemannian_cascade.pdf},
\path{fig_nonconvex_kl_pole.pdf}, and
\path{fig_nonconvex_clipped_kl_stationarity.pdf}.
\par\endgroup

\subsection{Fixed-step barrier}

All runs are one-dimensional, deterministic, and CPU-only; the potentials, their first two
derivatives, and the objective are closed form up to tabulated antiderivatives of the fixed
$C^\infty$ profiles, and the Gaussian averages $b(m,c)$, $A(m,c)$ use fixed Gauss--Hermite
quadrature, $32$ nodes for the bump train and $64$ for the two-cycle targets.

The bump-train grid instantiates \Cref{thm:sharp-disc} with profile constants
$T=1/8$, $Y=4$, $\LH=1$, $\gamma=1/2$, over
$\kappa\in\{8,16,\dots,1024\}$, in two families: the Appendix~\ref{app:bump} train built for
the certified step, and the same construction retuned so that the per-step mean gain
$s=\Delta t\,c_\kappa$ matches the tested arm. The two fixed arms are
$\Delta t=\gamma/\kappa$ and $\Delta t=\gamma$; a run stops at a $10^{-6}$ relative gap, and
the reported count $n_{1/2}$ is the first iterate with
$\DeltaE\leq\tfrac12\DeltaE_0$. Reported exponents are least-squares slopes of
$\log n_{1/2}$ against $\log\kappa$ over the upper half of the grid. A companion sweep
varies the order-one constant $\gamma\in\{1/8,\dots,4\}$; every $\gamma\leq1$ is monotone
and convergent at every $\kappa$ with exponent in $[0.88,0.92]$, and the fixed step first
fails at $\gamma=2$, independently of $\kappa$.

The two-cycle grid instantiates \Cref{thm:two-cycle} with the three-plateau profile of
\Cref{app:cycle} at $\LH=1$, over $\kappa\in\{64,128,256,512,1024\}$ and
$\gamma\in\{0.25,0.5,1,1.5\}$; the plateau length is obtained by a scalar root-find so that
$b(M,c)=rM$ holds to machine precision, and the realized diagnostics
($\min V''=1$, $\max V''=\kappa$, $\norm{V'''}_\infty=\LH$, $cA(M,c)=1$ and mean multiplier
$-1$ to twelve digits) are recorded per target. Cycle runs use a $2000$-step cap; the basin
study perturbs the initial mean by relative amounts in $\{0,10^{-12},10^{-6},
10^{-2},0.05,0.1,0.2,0.5\}$ with a $20000$-step cap. The Bures--Wasserstein contrast runs
the scalar forward--backward step on the same targets at
$\eta\in\{0.25,0.5,1,2,4\}/\beta$. The barrier panel of \Cref{fig:num-barrier} combines, at
$\kappa=256$, the retuned bump train over
$\Delta t\in\{2^{-11},\dots,2/\kappa\}$ with the two-cycle regime
$\Delta t\in(2/\kappa,1]$. Code and outputs:
\path{src/natural_gradient_fixed_step_barrier},
\path{outputs/natural_gradient_fixed_step_barrier}, with per-run CSVs
(\path{sharp_bump_*.csv}, \path{two_cycle_*.csv}) and the assertion suite
\path{tests/natural_gradient_fixed_step_barrier}.

\subsection{Provenance limits of the committed snapshot}

\begingroup\small\raggedright
The scientific parameters above are recoverable from the configurations,
source, and raw columns used by the selected figures, but some run-level
provenance files that the runners are designed to write are not committed:
\path{outputs/natural_gradient_covariance_bootstrap} lacks
\path{run_metadata.json};
\path{outputs/natural_gradient_stl_variance} lacks
\path{algorithm_results_long.csv}, \path{target_metadata.json},
\path{run_metadata.json}, \path{estimator_target_metadata.json}, and
\path{estimator_run_metadata.json}; and the Gaussian affine-mode runner does not
write a run-metadata file. The exact Python, NumPy, Torch, CUDA, and hardware
versions of those original runs therefore cannot be reconstructed from this
snapshot. The estimator-level STL CSV records the Torch/CUDA backend but not the
GPU model, the algorithm-level backend is not recorded in the surviving summary
files, and \path{requirements.txt} and \path{requirements-gpu.txt} give lower
dependency bounds rather than a locked environment. These omissions do not
change the selected data-to-figure mappings, since the two STL panels need only
the surviving estimator and tail CSVs and the covariance panels need only the
surviving long and transport-summary CSVs. The target-level metadata of
the covariance-bootstrap group is a deterministic function of the committed
configuration and has been regenerated from it
(\path{target_metadata.json}); the run-level metadata of the original runs is
not reconstructed. The asset builder \path{reports/make_report_assets.py}
treats the GPU-regenerable trajectory file \path{algorithm_results_long.csv}
and the algorithm run metadata as optional and builds every figure and table
used by this manuscript from the committed summary files, so both
\texttt{--only stl\_variance} and \texttt{--only covariance\_bootstrap} run to
completion from a clean checkout; the supplementary algorithm-gap panels, which
this manuscript does not use, require rerunning
\path{scripts/natural_gradient_stl_variance/run_algorithm_grid.py} on a CUDA
host. No missing value has been guessed here.
General WFR experiments under \path{outputs/wfr_gradient_flow} are not used in
any figure or claim in this manuscript.
\par\endgroup

% SOURCES:
% codex-draft/appendices/E-experimental-details.tex l.1-24 (preamble, objective,
%   Gauss-Hermite rescaling, asset-copy statement) -> opening paragraphs
% codex-draft/appendices/E-experimental-details.tex l.25-117 (covariance
%   bootstrap grids, Bures step, lineage) -> app subsection 1
% codex-draft/appendices/E-experimental-details.tex l.119-188 (discretization
%   grid, stability/accuracy definitions, lineage) -> app subsection 2
% codex-draft/appendices/E-experimental-details.tex l.190-288 (STL estimator and
%   algorithm grids, SHA256 cell seed, tail floors, lineage) -> app subsection 3
% codex-draft/appendices/E-experimental-details.tex l.290-379 (omega-tau modes,
%   Sobol designs, five initializations, lineage) -> app subsection 4
% codex-draft/appendices/E-experimental-details.tex l.381-457 (double-well
%   cascade, KL pole, clipped-KL stationarity, lineage) -> app subsection 5
% codex-draft/appendices/E-experimental-details.tex l.459-495 (provenance limits)
%   -> app subsection 6, condensed to one paragraph, content unchanged
% Cross-reference retargeting: fig:num-affine-modes -> fig:num-affine (label of
%   sections/06-numerics.tex l.163); other figure labels -> fig:num-three-stage,
%   fig:num-discretization, fig:num-stl, fig:num-nonconvex.
% Cross-references into the new namespace (sections/02-nonconvex-part.tex):
%   double-well display -> \eqref{eq:double-well} (l.28); Riemannian cascade map
%   -> \eqref{eq:cascade-map} and \Cref{prop:cascade-Riem} (l.41-46); KL
%   resolvent pole -> \eqref{eq:KL-cov-map} and \Cref{prop:pole-KL} (l.57-61);
%   split stationarity measure -> \eqref{eq:split-stationarity} and
%   \Cref{thm:clip-KL} (l.124, l.129); clipped update -> \eqref{upd:clip-KL}.
%   Modal rates -> \Cref{thm:aff-modes}. Scheme displays deduplicated against
%   \eqref{upd:Riem}, \eqref{upd:KL}; Radau reference solution -> \eqref{ODEs:NG}.
% Deduplication introducing two local labels: eq:app-diag-fixed-point (the
%   diagonal log-cosh fixed point, stated once and reused by the STL subsection)
%   and eq:app-affine-step (the omega-tau step, stated once and specialized to
%   the Gaussian calibration target).
% Notation retargeting per drafting-contract.md sec. 2: d -> N_\theta,
%   \pi -> \rho_\post, h -> \Delta t, x -> \theta, \lambda_0 -> \lambda_{0,\min},
%   energy gap -> \DeltaE; Gauss-Hermite nodes written (z_j,w_j) to free x.
